% theme: revolution_and_annual_motion
\MajorQuestion{公転と年周運動}
\SetRowSpaceQanda

\Instruction{次の問いに答えなさい。}
\QQ{天体が、ほかの天体のまわりを回ることを\NamedBlank{revolution}という。}{公転}
\QQ{太陽が真南を通るときの高度を\NamedBlank{culm_alt}という。}{南中高度}
\QQ{天球上の太陽の通り道を\NamedBlank{ecliptic}という。}{黄道}
\QQ{地球の北極と南極を結ぶ軸を\NamedBlank{earth_axis}という。}{地軸}
\QQ{\RefBlank{revolution}によって起こる天体の見かけの運動を\NamedBlank{annual_motion}という。}{年周運動}

\Instruction{\strut}
\QQ{地球が太陽のまわりを回る道筋を\NamedBlank{orbit}という。}{公転軌道}
\QQ{\RefBlank{ecliptic}付近にある12の星座を何というか。}{黄道12星座}
\QQ{1年のうち、北半球で\RefBlank{culm_alt}が最も高くなる日を何というか。}{夏至}
\QQ{太陽が星座の間を毎日少しずつ西から東へ動くように見える運動を\NamedBlank{solar_annual}という。}{太陽の年周運動}
\QQ{地球の\RefBlank{revolution}の向きは、自転の向きとどうなっているか。}{同じ向き}

\Instruction{\strut}
\QQ{地球が\RefBlank{revolution}するとき、地球が太陽を回る面を何というか。}{公転面}
\QQ{同じ時刻に見える星の位置が、1年周期で変わる見かけの動きを\NamedBlank{stellar_annual}という。}{星の年周運動}
\QQ{1年のうち、北半球で\RefBlank{culm_alt}が最も低くなる日を何というか。}{冬至}
\QQ{北の空では、星はどの星の近くを中心に1年で1回転するように見えるか。}{北極星}
\QQ{地球が1年で\RefBlank{revolution}する角度は何度か。}{360$^\circ$}

\Instruction{\strut}
\QQ{地球が\RefBlank{earth_axis}を傾けたまま\RefBlank{revolution}するために起こる変化を何というか。}{季節の変化}
\QQ{地球の\RefBlank{revolution}では、太陽のある側と反対側の空に何が見えるか。}{真夜中の星座}
\QQ{同じ時刻に見える星の位置は、1か月で約何度ずつ動くか。}{約30$^\circ$}
\QQ{太陽が真東から出て真西に沈み、昼と夜の長さがほぼ等しい日を2つ答えなさい。}{春分・秋分}
\QQ{星が同じ位置に見える時刻は、1か月で約何時間早くなるか。}{約2時間}
